% Chapter 7, Exercises 16--20.

\begin{exercise}[16]
  \label{ex:ch7-16}
  Let $K$ be a number field and let $S$ be a set of primes in $\gotho_K$. One
  says that $S$ has \emph{lower Dirichlet density}
  \[
    \underline{\delta}(S) = \liminf_{s \to 1^+} \frac{1}{\log\frac{1}{s-1}}
    \cdot \sum_{P \in S} \frac{1}{N(P)^s}
  \]
  provided that the indicated $\liminf$ exists. The \emph{upper Dirichlet
  density} $\overline{\delta}(S)$ is likewise defined via the $\limsup$.
  \begin{enumerate}[label=(\alph*)]
    \item Show that every $S$ has both a lower and an upper Dirichlet
      density. [\emph{Hint}: Use \RV{Proposition 7-31}, part~(ii), and also
      that a bounded sequence in $\RR$ has an infimum and a supremum.]
    \item Show that $S$ has a Dirichlet density (as defined previously) if
      and only if $\underline{\delta}(S) = \overline{\delta}(S)$.
    \item Let $T$ be any finite set of primes in $\gotho_K$. Show that
      $\underline{\delta}(S \cup T) = \underline{\delta}(S)$ and
      $\overline{\delta}(S \cup T) = \overline{\delta}(S)$. Conclude, in
      particular, that any finite set of primes has density zero.
    % The scan reads "N(P) = 1"; degree one means that N(P) = p is a
    % rational prime.
    \item Let $\mathscr{P}_K^1$ denote the set of degree-one primes of $K$;
      that is, prime ideals $P$ of $\gotho_K$ such that $N(P)$ is a rational
      prime. Show that for any set $S$ of primes,
      $\underline{\delta}(S \cap \mathscr{P}_K^1) = \underline{\delta}(S)$
      and $\overline{\delta}(S \cap \mathscr{P}_K^1) = \overline{\delta}(S)$.
      Conclude that the set of primes in $\gotho_K$ of degree greater than
      one has density zero.
  \end{enumerate}
\end{exercise}

\begin{exercise}[17]
  \label{ex:ch7-17}
  Let $E/K$ be a finite abelian extension of number fields with Galois group
  $G = \Gal(E/K)$. Let $U = N_{E/K}(C_E)$. Recall that $C_K/U$ is isomorphic
  to $G$ via class field theory. (See \RV{Theorem 6-6}.) For any character
  $\chi$ of $G$, let $\widetilde{\chi}$ denote the corresponding character of
  $C_K$ that is trivial on $U$. Keeping in mind that the extension $E/K$ is
  abelian, for any prime $P$ of $K$, we let
  \[
    \varphi_P = \left(\frac{E/K}{P}\right)
  \]
  denote the corresponding Frobenius element. Now set
  \[
    \chi(P) =
    \begin{cases}
      \chi(\varphi_P) & \text{if $P$ is unramified in $E$ or if $\chi$ is
      trivial}, \\
      0 & \text{otherwise},
    \end{cases}
  \]
  and put
  \[
    L(s, \chi) = \prod_P \bigl(1 - \chi(P) N(P)^{-s}\bigr)^{-1}.
  \]
  \begin{enumerate}[label=(\alph*)]
    \item Show that $L(s, \chi)$ converges absolutely in $\{\Re(s) > 1\}$ and
      further admits a meromorphic continuation to the whole $s$-plane, with
      no poles except possibly a simple one at $s = 1$, which occurs only
      when $\chi = 1$. Show, moreover, that there is a functional equation
      relating $s$ to $1 - s$. [\emph{Hint}: Use
      \RV{Theorems 6-6 and 7-16}.]
    % Note: with \chi(P) = 0 at every P ramified in E, (b) fails at a
    % ramified P that is not totally ramified; it holds if one sets
    % \chi(P) = \chi(\varphi_P) whenever \chi is trivial on the inertia group
    % at P (and up to finitely many Euler factors in any case, which
    % suffices for (c)).
    \item Show that
      \[
        \zeta_E(s) = \prod_{\chi \in \dual{G}} L(s, \chi).
      \]
      [\emph{Hint}: Write the left-hand side as an Euler product over the
      primes $P$ of $K$.]
    \item Show that
      \[
        \lim_{s \to 1^+} \frac{1}{\log\frac{1}{s-1}}
        \sum_{P \in \mathscr{P}_K^1} \frac{\chi(P)}{N(P)^s}
        =
        \begin{cases}
          1 & \text{if } \chi = 1, \\
          0 & \text{otherwise.}
        \end{cases}
      \]
  \end{enumerate}
\end{exercise}

\medskip
The following four problems lead the reader through a proof of a version of
the Tchebotarev density theorem, here reformulated in terms of Dirichlet
density. Recall that $\mathscr{P}_{K,f}$ denotes the set of finite places of
$K$, or, equivalently, the set of prime ideals of $\gotho_K$.
\begin{quote}
  \textsc{Theorem.} \emph{Let $E/K$ be a finite Galois extension of number
  fields with Galois group $G$, and let $C$ be a conjugacy class in $G$. Set}
  \[
    S_K(C) = \Bigl\{P \in \mathscr{P}_{K,f} :
    \left(\frac{E/K}{P}\right) = C\Bigr\}.
  \]
  \emph{Then $\delta(S_K(C))$ exists and equals
  $\operatorname{Card}(C)/\operatorname{Card}(G)$.}
\end{quote}
We shall also need the degree-one version of $S_K(C)$, which we define as
\[
  S_K^1(C) = S_K(C) \cap \mathscr{P}_K^1.
\]
The first problem is well known and deals with the abelian case. The next
three encapsulate the ideas from an elegant proof due to O.~Schreier.

\begin{exercise}[18]
  \label{ex:ch7-18}
  Let $E/K$ be a finite abelian extension, so that $C$ reduces to a singleton
  set $\{a\}$ for some $a \in G$. We recall from
  \RV{Chapter 3, Exercise 13}, the discrete form of the Fourier inversion
  formula: if $f$ is a complex-valued function on $G$, then
  \[
    f(g) = \frac{1}{\operatorname{Card}(G)}
    \sum_{\chi \in \dual{G}} \dual{f}(\chi)\chi(g)
  \]
  where
  \[
    \dual{f}(\chi) = \sum_{g \in G} f(g)\overline{\chi}(g).
  \]
  % Note: as printed, f_a is Card(G) times the characteristic function of
  % {a}, and the identity in (b) should read
  %   \sum_{P \in S_K(C)} N(P)^{-s}
  %     = \frac{1}{Card(G)} \sum_{g \in G} \sum_{P \in S_K(\{g\})} f_a(g) N(P)^{-s}
  %     = \frac{1}{Card(G)} \sum_{\chi} \overline{\chi}(a) \sum_P \chi(P) N(P)^{-s},
  % up to the finitely many primes ramified in E.
  \begin{enumerate}[label=(\alph*)]
    \item Given $a \in G$, define $f_a \colon G \to \CC$ by
      \[
        f_a(g) = \sum_{\chi \in \dual{G}} \chi(a^{-1}g).
      \]
      Show that $f_a$ is in fact the characteristic function of $\{a\}$.
    \item For all $g \in G$ show that
      \[
        \sum_{P \in S_K(C)} \frac{1}{N(P)^s}
        = \sum_{g \in G} \sum_{P \in S_K(C)} \frac{f_a(g)}{N(P)^s}
        = \frac{1}{\operatorname{Card}(G)}
        \sum_{\chi \in \dual{G}} \sum_P \frac{\chi(P)}{N(P)^s}.
      \]
    \item Show that $\delta(S_K(C))$ exists and equals
      $1/\operatorname{Card}(G) = \operatorname{Card}(C)/\operatorname{Card}(G)$.
      [\emph{Hint}: Use the previous exercise.]
  \end{enumerate}
\end{exercise}

\begin{exercise}[19]
  \label{ex:ch7-19}
  Suppose that $E/K$ is any finite normal extension of number fields. Let $G$
  be the associated Galois group, and let $C$ be a conjugacy class in $G$.
  Suppose further that $E'$ is an auxiliary finite abelian extension of $K$
  with Galois group $G'$ such that $E$ and $E'$ are linearly disjoint over
  $K$. Put $L = EE'$.
  \begin{enumerate}[label=(\alph*)]
    \item Show that $L/K$ is normal with Galois group
      $\mathcal{G} = G \times G'$.
    \item Choose any conjugacy class $C' = \{a\}$ in $G'$, and let $H$ be the
      subgroup generated by $CC'$ in $\mathcal{G}$. Put $F = L^H$, and
      consider the compositum $FE'$. We have the following diagram:
      \[
        \begin{tikzcd}[row sep=small, column sep=small]
          & & L = EE' \arrow[ddll, dash] \arrow[dd, dash, "H"']
            \arrow[dr, dash] & & \\
          & & & FE' \arrow[dl, dash] \arrow[dr, dash] & \\
          E \arrow[drr, dash, "G"'] & & F \arrow[d, dash] & &
            E' \arrow[dll, dash, "G'"] \\
          & & K & &
        \end{tikzcd}
      \]
      Show that $FE'$ is a normal extension of $F$ with Galois group given by
      \[
        \Gal(FE'/F) \cong H/(H \cap \langle C \rangle)
        \cong \langle C' \rangle \subseteq G'.
      \]
      Using the previous problem, conclude that
      \[
        \delta\Bigl(\Bigl\{Q \in \mathscr{P}_F^1 :
        \left(\frac{FE'/F}{Q}\right) = C'\Bigr\}\Bigr) = \frac{1}{f'}
      \]
      where $f' = \operatorname{Card}(\langle C' \rangle)$.
    \item For $P \in S_K^1(C)$, define
      $\alpha_P(F) = \operatorname{Card}(\{Q \in \mathscr{P}_F^1 : Q \mid P\})$.
      Show that
      \[
        \alpha_P(F) = \frac{1}{\operatorname{Card}(H)}
        \operatorname{Card}(\{\tau \in \mathcal{G} : \tau C \tau^{-1} \in H\}).
      \]
      [\emph{Hint}: Use \RV{Proposition 6-2}.]
    \item Put $f = \operatorname{Card}(C)$. Assuming that $f \mid f'$, show
      for $P \in S_K^1(C)$ that
      \[
        \left(\frac{L/K}{P}\right) \cap H = \varnothing.
      \]
  \end{enumerate}
\end{exercise}

\begin{exercise}[20]
  \label{ex:ch7-20}
  We continue in the context of the previous problem and always assume that
  $f \mid f'$.
  \begin{enumerate}[label=(\alph*)]
    \item Show that
      \[
        \lim_{s \to 1^+} \sum_{P \in S_K^1(C, C')} \frac{\alpha_P(F)}{N(P)^s}
        \sim \frac{1}{f'} \log\Bigl(\frac{1}{s-1}\Bigr)
      \]
      where
      \[
        S_K^1(C, C') = \Bigl\{P \in S_K^1(C) :
        \left(\frac{E'/K}{P}\right) = C'\Bigr\}.
      \]
    \item Show that
      \[
        \delta(S_K^1(C, C')) = \frac{f}{nn'}
      \]
      where $n$ and $n'$ denote $\operatorname{Card}(G)$ and
      $\operatorname{Card}(G')$, respectively. [\emph{Hint}: Show that
      $\alpha_P(F) = nn'/(ff')$.]
    \item Let
      \[
        n'_f = \operatorname{Card}(\{\sigma' \in G' : f \mid o(\sigma')\}).
      \]
      Show that
      \[
        \underline{\delta}(S_K(C)) \ge \frac{f n'_f}{nn'}
      \]
      where again $\underline{\delta}$ denotes lower Dirichlet density, as
      defined above. [\emph{Hint}: Any absolutely convergent summation
      indexed over $P \in S_K^1(C)$ can be rewritten as a double summation,
      first over $\sigma' \in G'$ and then over
      \[
        \Bigl\{P \in S_K^1(C) : \left(\frac{E'/K}{P}\right) = \sigma'\Bigr\}.
      \]
      Use this and appeal to part~(b).]
  \end{enumerate}
\end{exercise}
